\section{方法}

\begin{figure}[t]
\centering
\resizebox{\linewidth}{!}{%
\begin{tikzpicture}[
  font=\small,
  box/.style={draw=black!45, rounded corners=2pt, line width=0.55pt, align=center, minimum width=2.85cm, minimum height=0.82cm, fill=white},
  source/.style={box, fill=blue!5},
  evidence/.style={box, fill=cyan!5},
  knowledge/.style={box, fill=orange!9},
  model/.style={box, fill=green!8},
  output/.style={box, fill=purple!6},
  lane/.style={draw=black!18, rounded corners=4pt, fill=black!2},
  arrow/.style={-{Latex[length=2.1mm]}, line width=0.65pt, draw=black!70},
  softarrow/.style={-{Latex[length=1.8mm]}, line width=0.55pt, draw=black!45}
]
\node[source] (q) at (0,1.45) {\textbf{输入}\\医学问题 $q$};
\node[evidence] (retrieval) at (3.35,1.45) {\textbf{候选生成}\\BM25、稠密、fielded BM25、\\MedCPT 双编码器};
\node[evidence] (pool) at (6.95,1.45) {\textbf{Top-$M$ 证据池}\\验证集选择的\\加权 RRF};
\node[knowledge] (kg) at (3.35,-0.15) {\textbf{医学约束}\\实体、MeSH 层级、\\PrimeKG 辅助链接};
\node[knowledge] (hg) at (6.95,-0.15) {\textbf{局部超图}\\问题、段落、文档、\\实体与概念节点};
\node[model] (features) at (10.45,-0.15) {\textbf{可解释特征}\\语义分数、层级、\\扩散、中心性};
\node[model] (ltr) at (10.45,1.45) {\textbf{查询分组 LTR}\\LambdaMART 打分};
\node[output] (out) at (13.75,1.45) {\textbf{输出}\\排序证据 top-$k$\\与可选路径};

\begin{scope}[on background layer]
  \node[lane, fit=(q)(retrieval)(pool), inner sep=8pt] (laneA) {};
  \node[lane, fit=(kg)(hg)(features), inner sep=8pt] (laneB) {};
\end{scope}
\node[anchor=west, font=\footnotesize\bfseries, text=black!62] at ($(laneA.north west)+(0.12,-0.18)$) {检索层};
\node[anchor=west, font=\footnotesize\bfseries, text=black!62] at ($(laneB.north west)+(0.12,-0.18)$) {证据控制层};

\draw[arrow] (q) -- (retrieval);
\draw[arrow] (retrieval) -- (pool);
\draw[arrow] (pool) -- (ltr);
\draw[arrow] (ltr) -- (out);
\draw[softarrow] (kg) -- (hg);
\draw[softarrow] (pool) -- (hg);
\draw[softarrow] (q.south) |- (hg.west);
\draw[softarrow] (hg) -- (features);
\draw[softarrow] (features) -- (ltr);
\end{tikzpicture}}
\caption{KCH-MedRank 的论文风格方法总览。候选生成阶段提供高召回证据池；证据控制层将生物医学约束和局部超图结构转换为可解释特征；查询分组排序器输出用于下游 medical RAG 的最终证据顺序。}
\label{fig:method-overview}
\end{figure}


\subsection{第一阶段检索}

第一阶段在任何知识感知重排序之前先构建高召回候选池。我们使用 BM25 作为词法基线，稠密检索作为语义基线，并使用倒数排名融合（RRF）作为默认混合检索器。对于检索来源 $s$，记段落 $d$ 对问题 $q$ 的排名为 $r_s(q,d)$，融合分数为：
\[
  \mathrm{RRF}(q,d)=\sum_{s \in S}\frac{w_s}{k+r_s(q,d)},
\]
其中 $S$ 是检索来源集合，$w_s$ 是来源权重，所报告实验中 $k=60$。原始 Hybrid RRF 使用等权 BM25 和稠密检索。

增强候选生成路径的目标是提高重排序所需的 top-100 召回上限。该路径加入基于标题、正文和 MeSH 扩展字段的 fielded BM25，并加入 MedCPT dual-encoder 稠密检索作为额外生物医学语义来源。融合权重仅在验证集上选择。被选中的 w122 设置赋予 BM25 和稠密检索权重 1.0、fielded BM25 权重 1.2、MedCPT dual-encoder 权重 0.2。该设置根据验证集 Recall@100 选择，而不是根据测试集表现选择，因为重排序只能改善已经进入候选池的证据段落。

\subsection{知识约束局部超图}

对于每个问题及其 top-$M$ 候选，KCH-MedRank 构建局部超图 $H_q=(V_q,\mathcal{E}_q)$。节点集合包括问题节点、候选段落节点、文档节点、生物医学实体节点、MeSH 概念节点、可选 MeSH 树节点和可选关系节点。超边可以同时连接多个节点，因此能够表示普通成对相似度不易捕捉的高阶关系。

主要超边类型包括 question-entity、question-MeSH、question-passage candidate、passage-entity、passage-document、document-MeSH、shared-entity passages、entity-type concept、MeSH hierarchy，以及辅助 PrimeKG 关系边。用于消融的成对图变体和无医学知识超图变体来自同一候选池：成对图变体将多节点超边拆解为以锚点为中心的成对边，无医学知识变体只保留候选结构和排名分组结构。

局部设计避免在数百万生物医学段落上构建全局图。同时，该设计保持方法可解释：特征分析和案例研究可以报告哪些概念重叠、层级链接、共享候选簇或超图中心性信号影响了重排序。

\subsection{超图扩散与结构特征}

KCH-MedRank 将超图扩散作为可解释特征，而不是作为独立的重排序公式。种子分布在问题节点、问题实体节点和问题 MeSH 节点上赋予初始质量。每次迭代通过加权超边传播分数：
\[
  x^{(t+1)}=(1-\alpha)x^{(0)}+\alpha \cdot \mathrm{Propagate}(x^{(t)},\mathcal{E}_q),
\]
其中 $\alpha=0.85$，实现使用三轮扩散。段落级扩散分数随后在每个问题的候选集合内进行 min-max 归一化。我们还计算加权度中心性、局部节点数、局部超边数、共享实体边数、文档-MeSH 边数以及可选层级边数。

\subsection{生物医学知识特征}

重排序器使用若干具有医学动机的特征组。实体特征包括实体重叠数量、实体 Jaccard 相似度、问题实体覆盖率和段落实体数量。精确 MeSH 特征包括 MeSH 重叠、MeSH Jaccard 相似度、问题 MeSH 覆盖率和段落 MeSH 数量。层级感知 MeSH 特征包括精确描述符重叠、父级匹配、祖先匹配、兄弟节点匹配、树距离相似度、查询概念覆盖、段落层级覆盖、段落概念特异性，以及基于 MeSH 术语或父级概念的共享候选簇。

PrimeKG 派生特征仅作为辅助关系信号，包括关系数量、问题关系覆盖率和局部关系边数量。由于精确名称关系覆盖较稀疏，这些特征不被视为主要提升来源。

\subsection{学习排序}

主监督重排序器采用查询分组学习排序模型，使用 LightGBM/LambdaMART 风格的梯度提升决策树 \citep{ke2017lightgbm}。每个问题构成一个排序组，候选段落构成样本行，二值相关性标签来自金标准证据标注。特征向量结合检索分数、生物医学语义分数、实体特征、精确与层级感知 MeSH 特征、辅助关系特征和局部超图结构特征。

超参数在验证集上选择。搜索网格包括叶子数、学习率、估计器数量，以及可选的 blend weight，用于在模型分数和原始候选排名分数之间插值。选择准则优先考虑验证集 MRR@10，其次是 Recall@10 和 nDCG@10。完成验证选择后，最终模型在训练集加验证集上重新训练，并只在留出测试集上评估一次。

\begin{table}[t]
\centering
\small
\caption{增强 BioASQ 实验中 KCH-MedRank 的主要特征重要性。}
\label{tab:feature-importance-summary}
\begin{tabular}{lr}
\toprule
特征 & 重要性 \\
\midrule
生物医学语义分数 & 505 \\
Hybrid 分数 & 416 \\
Dense 分数 & 376 \\
BM25 排名分数 & 349 \\
生物医学语义排名分数 & 304 \\
段落实体数量 & 246 \\
Dense 排名分数 & 242 \\
BM25 分数 & 215 \\
共享 MeSH 术语簇大小 & 197 \\
共享 MeSH 父级簇大小 & 174 \\
超图度中心性 & 149 \\
段落 MeSH 特异性 & 135 \\
超图扩散分数 & 130 \\
\bottomrule
\end{tabular}
\end{table}



表~\ref{tab:feature-importance-summary} 的特征重要性显示，语义信号和第一阶段检索信号仍然处于核心位置，同时 MeSH 共享簇、段落概念特异性和超图中心性特征提供了可解释的生物医学与结构信息。这支持一个谨慎结论：KCH-MedRank 的重排序提升来自检索强度、生物医学语义和局部知识约束结构的结合，而不是依赖单一知识源。
